\section{The institutional gap}\label{sec:institutional}

\subsection{Ad-hoc intervention where an instrument is missing}

The UK's revealed response to trade-driven industrial distress is
firm-level rescue with improvised worker support attached. Two current
interventions define the pattern. At Port Talbot, the Tata Steel
transition is supported by a Transition Board fund of
\pounds\InstPortTalbotTotal\,million---\pounds\InstPortTalbotGov\,million
of UK Government money (an \pounds 80\,million package plus a
\pounds 22\,million addition announced in December 2025) and
\pounds 20\,million from Tata Steel---funding local support for affected
workers and businesses \citep{ukgov2025porttalbot}. At Scunthorpe, the
National Audit Office's March 2026 investigation into the British Steel
intervention found \pounds\InstScunthorpeSpend\,million of departmental
spending between April 2025 and January 2026, running at roughly
\pounds\InstScunthorpeDaily\,million a day, with no set budget, no
repayment schedule and no end date \citep{nao2026}. The Public Accounts
Committee took oral evidence on \InstPACDate\ from the steel unions and
the local authority on what the intervention and possible nationalisation
mean for workers \citep{pac2026}. Taken together, government money
committed around the two steelworks is roughly
\pounds\InstAdHocSteelTotal\,million and rising---spent without a
designed worker-side instrument, without stated eligibility rules a
displaced worker elsewhere could invoke, and without the evaluation
apparatus a statutory scheme would carry.

The only standing machinery for redundant workers, the Jobcentre Plus
Rapid Response Service, provides job-search support, training referrals
and grants for specific costs such as travel and tools; it pays no income
replacement \citep{dwprrs}. Income replacement for a displaced worker is
the ordinary benefit system: contributory New Style JSA at
\pounds\TaaJsaWeekly\ a week for 26 weeks, then means-tested
Universal Credit---gated by the capital test into which statutory
redundancy pay itself counts.

\subsection{A live reform that is almost the missing instrument}

The March 2025 \emph{Pathways to Work} Green Paper proposed replacing New
Style JSA and New Style ESA with a single time-limited contributory
\emph{Unemployment Insurance}, paid at the higher ESA rate
(\pounds\InstUIWeeklyRate\ a week at 2025--26 rates) with a duration
left to consultation---six to twelve months was the range
consulted on \citep{dwp2025pathways}. The October 2025 consultation
response summarised submissions without committing to a rate, duration
or timetable, and no legislation has been introduced
\citep{dwp2025response}; the proposal stands, unlegislated, as of August
2026. For this paper's purpose its status is exactly right: the
government has already designed most of Arm~B. What it has not done is
ask what its own proposal would deliver to trade-displaced workers
specifically, or what extending it toward the two-year horizons of the US
programmes would cost---the incremental questions
Section~\ref{sec:results} prices.

\subsection{The asymmetry in the trade strategy}

The June 2025 Trade Strategy commits to ``significantly upgrading our
trade defence toolkit'' and overhauling the trade-remedies system, and
the steel safeguard's replacement from July 2026 tightened import
protection further \citep{dbt2025strategy}. The strategy contains no
income-support, compensation or adjustment-assistance instrument for
workers displaced by trade. The asymmetry is complete: firms facing
import competition have a statutory, rules-based remedies system with an
independent authority; workers facing the same shock have discretionary,
place-specific packages when their employer is large enough to rescue.
The four instruments of Section~\ref{sec:design} are candidate
corrections to that asymmetry, sized so the comparison with the ad-hoc
spending is direct.
